%% kapa: exam=1 round=33 year=2022 period=1 form=A subject=real_estate subject_ko=부동산학원론 seq=16 number=96
%% source: Q-Net 감정평가사 기출문제 (cst00309) artlSeq=5209997 file=raw/1008/5209997/2022년도 제33회 감정평가사 1차시험 1교시(민법,경제학,부동산학).pdf
%% conversion: kordoc 4.24.51 -> tools/gen1.py (auto); review=proofread-sonnet; flags=none
\begin{question}{96}
부동산 투자분석기법에 관한 설명으로 옳은 것은?
\begin{choices}
  \item 투자규모가 상이한 투자안에서 수익성지수(PI)가 큰 투자안이 순현재가치(NPV)도 크다.
  \item 서로 다른 투자안 A, B를 결합한 새로운 투자안의 내부수익률(IRR)은 A의 내부수익률과 B의 내부수익률을 합한 값이다.
  \item 순현재가치법과 수익성지수법에서는 화폐의 시간가치를 고려하지 않는다.
  \item 투자안마다 단일의 내부수익률만 대응된다.
  \item 수익성지수가 1보다 크면 순현재가치는 0보다 크다.
\end{choices}
\end{question}
\answer{5}
